\section{Chameleon: convertir todas las modalidades en tokens discretos}

La primera ruta persigue la unificación de interfaces en su máxima expresión: si tanto el texto como la imagen pueden convertirse en IDs discretos, se puede hacer uso directo del modelado de lenguaje causal. El significado histórico de Chameleon radica en demostrar que la fusión temprana multimodal (mixed-modal early fusion) es viable durante el preentrenamiento a gran escala; además, deja clara las desventajas de una interfaz unificada: la compresión mediante codebook provoca la pérdida de semántica visual fina, y la eficiencia predictiva discreta para la generación de imágenes no necesariamente supera a la difusión latente continua.

\subsection{Supuestos fuertes de la fusión temprana}

La primera opción dejada pendiente en el capítulo anterior es la representación. En esta sección comenzamos con la solución más radical: discretizar la imagen y completar la fusión antes de que los datos ingresen al Transformer. Al analizar la figura, primero observe dónde se ubica el tokenizador y luego considere qué ventajas de entrenamiento ofrece la secuencia unificada.

\lecturefigure{slide-022.jpg}{Modelo de fusión temprana (early-fusion) de Chameleon: convertir todas las modalidades en tokens discretos}{Diapositiva 22 del deck oficial; arXiv de Chameleon 2405.09818v2}

\begin{knowledgebox}{Lectura de la figura: la fusión temprana ocurre antes que el modelo reciba la entrada}
En la imagen, primero el tokenizador convierte la imagen en tokens visuales discretos; luego, estos se intercalan con los tokens de texto para ingresar al mismo Transformer. Lo que denominamos fusión temprana (early fusion) implica formar una secuencia unificada antes de las capas profundas del modelo, en lugar de ejecutar por separado el codificador visual y el modelo de lenguaje y después concatenar sus logits. Su ventaja es que todas las posiciones pueden participar en el mismo contexto causal; su riesgo es que la botella de información (information bottleneck) del tokenizador visual se traslada a todas las tareas posteriores.
\end{knowledgebox}

Sea $z_e(x)$ el latente continuo generado por el encoder y $\{e_k\}_{k=1}^{K}$ el codebook. La asignación VQ (Vector Quantization) se escribe como:
$$
k^*(x)=\arg\min_{k\in\{1,\ldots,K\}}\left\|z_e(x)-e_k\right\|_2^2,
\qquad z_q(x)=e_{k^*(x)}.
$$
$K$ es el tamaño del codebook y $z_q$ es el latente cuantizado. A mayor compresión, la secuencia será más corta; sin embargo, si varios parches con detalles distintos caen en el mismo código, el modelo no podrá recuperar la información perdida a partir del ID del token.

\lecturefigure{slide-023.jpg}{VQ-VAE mapea los parches de imagen continuos al codebook discreto}{Diapositiva 23 del deck oficial; VQ-VAE-2 arXiv 1906.00446v1; clase 00:12:55--00:14:18}

\begin{knowledgebox}{Lectura de la figura: el VQ-VAE asume simultáneamente compresión y construcción del vocabulario}
El encoder izquierdo comprime los píxeles en una cuadrícula latente; el cuantizador central selecciona para cada posición el vector del codebook más cercano; el decoder derecho intenta reconstruir la imagen. Para Chameleon, el resultado más importante no es la propia imagen reconstruida, sino la secuencia de índices del codebook. Esta secuencia puede ingresar a la entropía cruzada (cross-entropy) como lo haría un texto, pero la calidad de reconstrucción y la fidelidad semántica están conjuntamente limitadas por la capacidad del codebook, la tasa de submuestreo y el objetivo de entrenamiento.
\end{knowledgebox}

\begin{lstlisting}[caption={Pseudocódigo mínimo para tokenización de imágenes VQ}]
def vq_tokenize(image, encoder, codebook):
    latent = encoder(image)                 # [H, W, D]
    distance = pairwise_l2(latent, codebook)
    token_ids = distance.argmin(axis=-1)    # [H, W]
    return token_ids.flatten()
\end{lstlisting}

\teachervoice{00:12:27--00:12:55, el docente menciona que considerar "el mundo multimodal como una secuencia de tokens discretos" es una visión muy poderosa. Lo poderoso radica en la transferencia directa del stack de entrenamiento de LM; sin embargo, la descripción "powerful" se refiere a la uniformidad del modelado y no implica que la representación esté libre de pérdidas.}

\subsection{¿Por qué una secuencia puede generar cualquier orden de texto-imagen?}

Cuando los códigos de texto e imagen se insertan en la misma secuencia causal, el modelo puede aprender:
$$
p(t_1,\ldots,t_i,v_1,\ldots,v_j,t_{i+1},\ldots),
$$
por lo que puede generar tanto de texto a imagen como de imagen a texto, además de producir alternativamente múltiples segmentos de texto e imágenes. La clave no reside en añadir un nuevo "modo especial" de imagen, sino en que los datos de entrenamiento ya contienen un orden intercalado (interleaved ordering); el modelo aprende a continuar prediciendo correctamente el espacio objetivo después del límite de modalidad (modality boundary).

\lecturefigure{slide-024.jpg}{Generación interleaved de texto-imagen de orden arbitrario en Chameleon}{Diapositiva 24 del deck oficial; clase 00:14:23--00:14:50}

\begin{knowledgebox}{Lectura de la figura: el orden arbitrario proviene de la distribución de entrenamiento, no de una combinación libre incondicional}
La imagen muestra que puede aparecer primero la figura y luego el texto, o viceversa, o alternarse; esto indica que el modelo posee capacidad de continuación multimodal (mixed-modal continuation). No obstante, sigue sujeto a limitaciones como el tamaño de la ventana de contexto, los patrones de interleaving vistos durante el entrenamiento y la fidelidad del tokenizador. Si los datos contienen casi exclusivamente pares de descripciones (caption pair), el modelo no aprenderá automáticamente una compleja orquestación de historias visuales-textuales solo por tener una arquitectura unificada.
\end{knowledgebox}

\lecturefigure{slide-025.jpg}{Capacidad multitask multimodal de Chameleon}{Diapositiva 25 del deck oficial; clase 00:14:50--00:15:26}

\begin{knowledgebox}{Lectura de la figura: la diversidad de tareas es la prueba real del valor de la fusión temprana}
La imagen presenta simultáneamente descripción (captioning), preguntas visuales (visual question answering), generación de imágenes y razonamiento multimodal. Primero compare las modalidades de entrada/salida, luego evalúe los criterios de evaluación: las tareas de comprensión suelen devolver texto, por lo que se usan métricas de exactitud o lingüísticas; las tareas de generación devuelven imágenes, requiriendo evaluación perceptual o humana. Listarlas en una sola página no implica que sus puntajes sean directamente sumables, ni que una única representación sea óptima para todas las tareas.
\end{knowledgebox}

\subsection{Dos tipos de cuellos de botella en la interfaz unificada de tokens}

Los resultados de Chameleon demuestran que es "viable", pero la clase posteriormente enfatiza dos límites. Primero, los códigos de imagen discretos suelen ser inferiores para la comprensión fina (fine-grained understanding) frente a embeddings semánticos continuos como SigLIP; segundo, el lado de generación mediante difusión/flow matching avanza más rápido en calidad y eficiencia. Surge así una contradicción central: es difícil que una única representación de imagen sirva simultáneamente para la comprensión semántica y la generación de alta fidelidad.

\lecturefigure{slide-026.jpg}{Contribuciones de Chameleon, limitaciones de la representación y posición tecnológica en 2026}{Diapositiva 26 del deck oficial; clase 00:15:32--00:16:30}

\begin{knowledgebox}{Lectura de la figura: el segundo punto del resumen es más importante que el primero}
El primer punto confirma que la fusión temprana desde cero puede mantener el rendimiento en solo texto y generar capacidad multimodal. El segundo punto restringe la conclusión: los tokens VQ sufren pérdida de semántica visual fina, y el lado de generación ha sido superado por la difusión latente continua. La nota "y muchos otros" en la esquina inferior derecha indica que modelos posteriores adoptan ampliamente codificadores de entrada continuos o representaciones híbridas; no se está afirmando que Chameleon haya sido completamente obsoleto.
\end{knowledgebox}

\begin{warningbox}{Conclusiones que no se pueden extraer de Chameleon}
No se puede concluir que todas las modalidades deban ser discretas; no se puede afirmar que un único vocabulario elimine la brecha entre modalidades (modality gap); no se puede declarar mediante muestras cualitativas que el razonamiento y la generación han experimentado transferencia positiva; tampoco se debe interpretar el mantenimiento del rendimiento en solo texto como alcanzar el estado del arte en todas las tareas visuales. La evidencia central de Chameleon es la viabilidad de la fusión temprana, no la optimalidad universal de la representación.
\end{warningbox}

\subsection{Resumen de la sección}

Chameleon transfiere el modelado de lenguaje de la manera más radical: todas las entradas y salidas se convierten en tokens discretos y se utiliza una entropía cruzada unificada. Obtiene interfaces elegantes de interleaving y multitask, y expone claramente la botella de compresión de la información visual. La siguiente ruta conserva un Transformer compartido, pero renuncia al requisito de que "todas las salidas deban ser IDs discretos".


